\subsection{জেনে রাখো: Non-reversible transformation}
এটি Encryption/Decryption পদ্ধতি নয়; একটি প্রদত্ত নিয়মে অক্ষর বদলানোর হিসাব। একই ফল থেকে মূল অক্ষর একার্থে উদ্ধার করা যায় না—এই সীমাবদ্ধতা যাচাই করাই উদাহরণটির উদ্দেশ্য। ধরি, $A=1, B=2,\ldots,Z=26$ এবং অবস্থানকে 8 দিয়ে গুণ করে 26 দিয়ে ভাগশেষ নিতে হবে; ভাগশেষ 0 হলে Z ধরা হবে। এই নিয়মে CAESAR-এর হিসাব নিচে দেওয়া হলো।
\begin{chSixSideBySide}[.46]
\small
\begin{tabular}{|l|l|l|l|}
\hline

\chSixHeader
অক্ষর
&
অবস্থান
&
গুণফল ও ভাগশেষ
&
ফল
\\ \hline

C
&
3
&
$24\bmod26=24$
&
X
\\ \hline

A
&
1
&
$8\bmod26=8$
&
H
\\ \hline

E
&
5
&
$40\bmod26=14$
&
N
\\ \hline

S
&
19
&
$152\bmod26=22$
&
V
\\ \hline

A
&
1
&
$8\bmod26=8$
&
H
\\ \hline

R
&
18
&
$144\bmod26=14$
&
N
\\ \hline

\end{tabular}
\chSixNextColumn
\textbf{হিসাব পড়ার নিয়ম:} C-এর অবস্থান 3; $3\times8=24$, তাই ফল X। E-এর জন্য $5\times8=40$ এবং $40\bmod26=14$, তাই ফল N। প্রতিটি অক্ষরে একই নিয়ম প্রয়োগ করে পাওয়া যায় \texttt{CAESAR}$\rightarrow$\texttt{XHNVHN}।

এখানে E ও R দুটির ফলই N। কারণ $\gcd(8,26)=2\ne1$; 8 ও 26 পরস্পর সহমৌলিক নয়। ফলে এই নিয়মে শুধু রূপান্তরিত ফল দেখে প্রতিটি মূল অক্ষর একার্থে ফেরত পাওয়া যায় না।
\end{chSixSideBySide}
\begin{keypointbox}{নিয়ম বুঝে হিসাব করবে}
Caesar-এর Shift Rule এবং এই গুণ করার নিয়ম আলাদা। এটি আধুনিক নিরাপদ অ্যালগরিদমের উদাহরণ নয়; উদ্দেশ্য হলো প্রশ্নের দেওয়া নিয়ম থেকে ধাপে ধাপে ফল নির্ণয় এবং তার সীমাবদ্ধতা বোঝা।
\end{keypointbox}

